\section{Study of fibres of a universally open morphism}
\label{section:IV.15}

\subsection{Multiplicities of fibres of a universally open morphism}
\label{subsection:IV.15.1}

\begin{proposition}[15.1.1]
\label{IV.15.1.1}
\oldpage[IV-3]{223}
Let $Y$ be an irreducible locally Noetherian prescheme, of generic point $\eta$, $f:X\to Y$ a morphism locally of finite type, $y$ a point of $Y$, $\sh{F}$ a coherent $\sh{O}_X$-Module; set $\sh{F}_y=\sh{F}\otimes_{\sh{O}_Y}k(y)$ (sheaf of modules on the fibre $f^{-1}(y)$).
Let $z$ be the generic point of an irreducible component of $\Supp(\sh{F})$, and let $X_i$ ($1\leqslant i\leqslant n$) be the closed reduced sub-preschemes of $X$ whose underlying spaces are the irreducible components of $\Supp(\sh{F})$ containing $z$, and which are such that the restriction $X_i\to Y$ of $f$ is a universally open morphism at the point $z$; let $z_i$ be the generic point of $X_i$ ($1\leqslant i\leqslant n$).
If $\lambda_z(\sh{F}_y)$ (resp.\ $\lambda_i(\sh{F}_\eta)$) denotes the geometric multiplicity of $\sh{F}_y$ at a point $x\in f^{-1}(y)$ (resp.\ that of $\sh{F}_\eta$ at a point $t\in f^{-1}(\eta)$) \sref[IV]{IV.4.7.5}, then
\[
\tag{15.1.1.1}
  \lambda_z(\sh{F}_y) \geqslant \sum_{i=1}^{n} \lambda_{z_i}(\sh{F}_\eta).
\]
If moreover the ring $\sh{O}_y$ is regular, then
\[
\tag{15.1.1.2}
  \operatorname{long}((\sh{F}_y)_z) \geqslant \sum_{i=1}^{n} \operatorname{long}((\sh{F}_\eta)_{z_i}).
\]
\end{proposition}

\begin{proof}
The fibres $f^{-1}(y')$ of $f$ (for all $y'\in Y$) do not change when we replace $f$ by $f_\mathrm{red}$, so we may suppose that $X$ and $Y$ are reduced (2.4.3,~(vi)), hence $Y$ integral; we set $K=k(\eta)$.
We may moreover replace $f$ by the restriction $\Supp(\sh{F})\to Y$ of $f$ to the closed sub-prescheme having $\Supp(\sh{F})$ as underlying space, and thus reduce to the case where $X=\Supp(\sh{F})$.
Finally, if $y=\eta$, there is only one irreducible component $X_i$ of $X$ containing $z$ (since $z_i=z$), and (15.1.1.1) and (15.1.1.2) are trivial.
We may thus restrict to the case $y\neq\eta$; it then follows that the $z_i$ belong to $f^{-1}(\eta)$, and so the second members of (15.1.1.1) and (15.1.1.2) are well-defined.
Let $Z_i$ be the closed reduced sub-prescheme of $f^{-1}(\eta)$ having $X_i\cap f^{-1}(\eta)$ as underlying space.
Since for each $i$, the $K'$-prescheme $(Z_i\otimes_K K')_\mathrm{red}$ is geometrically reduced over $K'$, there exists a finite radiciel extension $K'$ of $K$ such that the projection $f^{-1}(\eta)\otimes_K K'\to f^{-1}(\eta)$ is finite, dominant, and radiciel (\sref[I]{I.3.5.7}), hence a homeomorphism (2.4.5).
So $z_i'$ is the unique point of $f^{-1}(\eta)\otimes_K K'$ above $z_i$.
It follows from (4.7.9) and (4.7.5) that
\[
\tag{15.1.1.3}
  \lambda_{z_i}(\sh{F}_\eta)=\lambda_{z_i'}(\sh{F}_\eta\otimes_K K')_{z_i'}=\operatorname{long}((\sh{F}_\eta\otimes_K K')_{z_i'}).
\]

Let $V$ be a discrete valuation ring with field of fractions $K'$, dominating $\sh{O}_y$ (\textbf{II},~7.1.7); set $Y'=\Spec(V)$, $X'=X\times_Y Y'$, $\sh{F}'=\sh{F}\otimes_{\sh{O}_Y} Y'$, $f'=f_{(Y')}:X'\to Y'$; if $\eta'$ is the generic point of $Y'$ and $y'$ its closed point, then $g(y')=y$, $g(\eta')=\eta$; the morphism $\Spec(k(\eta'))\to\Spec(k(\eta))$ being faithfully flat, hence (2.2.13) and (2.3.4), there is a generic point $z'$ of an irreducible component of $f'^{-1}(y')$ such that $g'(z')=z$.
Par construction, $k(\eta')=K'$, hence $f'^{-1}(\eta')=f^{-1}(\eta)\otimes_K K'$; on the other hand, if we set $X_i'=X_i\times_Y Y'$, the restriction $X_i'\to Y'$ of $f'$ is a universally open morphism at the point $z'$, hence (1.10.3) there exists a generisation $z_i''$ of $z'$, belonging to an irreducible component of $X_i'\cap f'^{-1}(\eta')$, so we may suppose $z_i''=z_i'$.
By (4.7.9) and (4.7.5),
\[
\tag{15.1.1.4}
  \lambda_z(\sh{F}_y) = \lambda_{z'}(\sh{F}_{y'}) \geqslant \operatorname{long}((\sh{F}_{y'}')_z)
\]
and it follows from this inequality and (15.1.1.3) that to establish (15.1.1.1), it suffices to show the inequality
\[
\tag{15.1.1.5}
  \operatorname{long}((\sh{F}_{y'}')_z) \geqslant \sum_{i=1}^{n} \operatorname{long}((\sh{F}_{\eta'}')_{z_i'}).
\]

Now consider the second inequality (15.1.1.2); we use the following lemma:

\begin{lemma}[15.1.1.6]
\label{IV.15.1.1.6}
Let $A$ be a local regular ring that is not a field, $k$ its residue field, $K$ its field of fractions.
There exists a discrete valuation ring $V$ dominating $A$, having $K$ for field of fractions, and whose residue field is a pure transcendental extension of $k$.
\end{lemma}

Set $S=\Spec(A)$, and consider the $S$-scheme $P$ obtained by blowing up the closed point $s$ of $S$ (\textbf{II},~8.1.3); if $\mathfrak{m}$ is the maximal ideal of $A$, then by definition $P=\operatorname{Proj}(B)$, where $B$ is the graded $A$-algebra $\bigoplus_{n\geqslant 0}\mathfrak{m}^n$.
If $h:P\to S$ is the structural morphism, the fibre $h^{-1}(s)$ is isomorphic to $\operatorname{Proj}(B\otimes_A k)$ (\textbf{II},~2.8.10), and by definition $B\otimes_A k$ is the graded algebra $\bigoplus_{n\geqslant 0}\mathfrak{m}^n/\mathfrak{m}^{n+1}$; but since $A$ is regular, this algebra is isomorphic to a polynomial algebra $k[t_1,\ldots,t_e]$ ($t_i$ indeterminates) ($\mathbf{0}$,~17.1.1) and hence $h^{-1}(s)\cong\mathbf{P}_k^{e-1}$.
Let $\zeta$ be the generic point of $h^{-1}(s)$; if $t_i'$ ($1\leqslant i\leqslant e$) are the classes of generators of $\mathfrak{m}$ modulo $\mathfrak{m}^2$, the local ring $B_{(t_i')}$ is an affine neighbourhood of $\zeta$ in $P$, and we have $k(\zeta)=k(t_1,\ldots,t_{e-1})$; on the other hand, since $A$ is integrally closed, $B_{(t_i')}$ is also integrally closed, and hence the local ring $\sh{O}_{P,\zeta}$ is regular.
Finally, as $A=\sh{O}_s$, $\dim(\sh{O}_{P,\zeta})=\dim(\sh{O}_s)-(e-1)$, since $h$ is a birational morphism, that $\dim(h^{-1}(s))=e-1$, and the local ring of $P$ at its generic point $\zeta$ is a discrete valuation ring of dimension $1$, and $\sh{O}_{P,\zeta}=V$ is a discrete valuation ring that answers the question (\textbf{II},~7.1.6).

This lemma being established, we can resume the reasoning for the inequality (15.1.1.1) with $Y'=\Spec(V)$; this time $k(y')$ is a \emph{separable} extension of $k(y)$, and by (4.7.9), $\operatorname{long}((\sh{F}_y)_z)=\operatorname{long}((\sh{F}_{y'}')_{z'})$; as moreover $f'^{-1}(\eta')=f^{-1}(\eta)$ and $\sh{F}_{\eta'}'=\sh{F}_\eta$, we are again reduced to proving (15.1.1.5).
But if $t$ is a uniformiser of $\sh{O}_{y'}$, $f'^{-1}(y')$ is the closed sub-prescheme of $X'$ defined by the ideal $t\sh{F}'$; on the other hand, (15.1.1.5) is then just a special case of (3.4.1.1).
\end{proof}

\begin{corollary}[15.1.2]
\label{IV.15.1.2}
Under the hypotheses of \sref{IV.15.1.1}, suppose moreover that $\lambda_z(\sh{F}_y)=1$ (resp.\ that $\sh{O}_y$ is regular and $\operatorname{long}((\sh{F}_y)_z)=1$).
Then there is at most one irreducible component $X_i$ of $\Supp(\sh{F})$ containing $z$ and such that the restriction $X_i\to Y$ of $f$ is universally open at the point $z$.
Moreover, if $z_i$ is the generic point of this component, then $\lambda_{z_i}(\sh{F}_\eta)=1$ (resp.\ $\operatorname{long}((\sh{F}_\eta)_{z_i})=1$).
\end{corollary}

This follows immediately from \sref{IV.15.1.1}, noting that one necessarily has $(\sh{F}_\eta)_{z_i}\neq 0$ by definition of $X_i$ (\sref[I]{I.9.1.13}).

\begin{corollary}[15.1.3]
\label{IV.15.1.3}
Let $Y$ be an irreducible locally Noetherian prescheme of generic point $\eta$, $f:X\to Y$ a morphism locally of finite type, $\sh{F}$ a coherent $\sh{O}_X$-Module of support $X$, $x$ a point of $X$, $y=f(x)$.
Suppose that: $1\textsuperscript{o}$ $x$ belongs to a single irreducible component $Z$ of $f^{-1}(y)$; $2\textsuperscript{o}$ $\sh{F}_y=\sh{F}\otimes_{\sh{O}_Y}k(y)$ is geometrically reduced on $k(y)$, i.e.\ if $z$ is the generic point of $Z$, $\lambda_z(\sh{F}_y)=1$ (resp.\ $\sh{O}_y$ is regular and $\operatorname{long}((\sh{F}_y)_z)=1$).
Then there is at most one irreducible component $X'$ of $X$ containing $x$, such that $\dim_x(X'\cap f^{-1}(y))=\dim_x(f^{-1}(y))$ and such that the restriction $X'\to Y$ of $f$ is universally open at the generic points of the irreducible components of $f^{-1}(y)$ containing $x$.
Moreover, if such a component $X'$ exists, and if $z'$ is its generic point, then $\lambda_{z'}(\sh{F}_\eta)=1$ (resp.\ $\operatorname{long}((\sh{F}_\eta)_{z'})=1$).
\end{corollary}

Indeed, an irreducible component $X'$ of $X$ containing $x$ and such that $\dim_x(X'\cap f^{-1}(y))=\dim_x(f^{-1}(y))$ necessarily contains $Z$, since an irreducible component of $f^{-1}(y)$ containing $x$ must be of the same dimension as the sole irreducible component of $f^{-1}(y)$ containing $x$.
One may then apply \sref{IV.15.1.2} to $z$.

\begin{remarks}[15.1.4]
\label{IV.15.1.4}
\begin{enumerate}[label=(\roman*)]
  \item In the statement of \sref{IV.15.1.1}, one cannot replace ``universally open'' by ``equidimensional'', as the example (14.4.10,~(ii)) where one takes $\sh{F}=\sh{O}_X$ shows; the fibres of $f$ are then reduced Artinian schemes, so $\lambda_x(\sh{F}_y)=1$, but there are two irreducible components $X_1$, $X_2$ of $X$ containing $x'$, and the second member of (15.1.1.1) is therefore equal to $2$, even though the restriction of $f$ to each of the components $X_1$, $X_2$ is an equidimensional morphism at every point.
  This same example also shows that in \sref{IV.15.1.2} and \sref{IV.15.1.3}, one cannot replace the hypothesis ``universally open'' by ``equidimensional''.
  \item It seems likely that for the validity of the inequality (15.1.1.2), one cannot suppress the hypothesis that $\sh{O}_y$ is \emph{regular}.
\end{enumerate}
\end{remarks}

\subsection{Flatness of universally open morphisms with geometrically reduced fibres}
\label{subsection:IV.15.2}

\begin{env}[15.2.1]
\label{IV.15.2.1}
\oldpage[IV-3]{226}
We have seen (2.4.6) that a morphism \emph{flat} and locally of finite presentation is universally open.
One has a partial converse of this result, subject to supplementary hypotheses:
\end{env}

\begin{theorem}[15.2.2]
\label{IV.15.2.2}
Let $Y$ be a locally Noetherian prescheme, $f:X\to Y$ a morphism locally of finite type, $\sh{F}$ a coherent $\sh{O}_X$-Module, $x$ a point of $X$, $y=f(x)$.
Suppose the following conditions are satisfied:
\begin{enumerate}[label=(\roman*)]
  \item $\Supp(\sh{F})=X$.
  \item $f$ is universally open at the generic points of the irreducible components of $f^{-1}(y)$ containing $x$.
  \item $\sh{F}_y=\sh{F}\otimes_{\sh{O}_Y}k(y)$ is geometrically reduced at the point $x$ (4.6.22).
  \item The ring $\sh{O}_y$ is reduced.
\end{enumerate}
Then $\sh{F}$ is $f$-flat at the point $x$.
\end{theorem}

\begin{proof}
Since $\sh{O}_y$ is reduced and Noetherian, we shall apply the valuative criterion of flatness (11.8.1).
Consider therefore a local homomorphism $\sh{O}_y\to A'$, where $A'$ is a discrete valuation ring; set $Y'=\Spec(A')$, $X'=X\times_Y Y'$, $f'=f_{(Y')}:X'\to Y'$; let $x'$ be a point of $X'$ whose projections in $X$ and $Y'$ are respectively $x$ and the closed point of $Y'$; if $\sh{F}'=\sh{F}\otimes_{\sh{O}_Y} Y'$, then $\Supp(\sh{F}')=X'$ by (\sref[I]{I.9.1.13}); every generic point $z'$ of an irreducible component of $f'^{-1}(y')$ containing $x'$ has for projection in $f^{-1}(y)$ a generic point $z$ of an irreducible component of $f^{-1}(y)$ containing $x$, by (4.2.6); hence $f'$ is universally open at $z'$.
Finally, it follows from (4.7.11) that $\sh{F}_{y'}'$ is geometrically reduced at the point $x'$.
One thus sees that one is reduced to showing the following:
\end{proof}

\begin{lemma}[15.2.2.1]
\label{IV.15.2.2.1}
Let $Y=\Spec(A)$ be the spectrum of a discrete valuation ring, $y$ its closed point, $f:X\to Y$ a morphism locally of finite type, $x$ a point of $f^{-1}(y)$, $\sh{F}$ an $\sh{O}_X$-Module coherent.
Suppose the following conditions are satisfied:
\begin{enumerate}[label=(\roman*)]
  \item $\Supp(\sh{F})=X$.
  \item $f$ is open at the generic points of the irreducible components of $f^{-1}(y)$ containing $x$.
  \item $\sh{F}_y=\sh{F}\otimes_{\sh{O}_y}k(y)$ is reduced at the point $x$ (3.2.2).
\end{enumerate}
Then $\sh{F}$ is $f$-flat at the point $x$.
\end{lemma}

\begin{proof}
Let $t$ be a uniformiser of $A$; set $B=\sh{O}_x$ and $\sh{F}_x=M$; condition (i) implies that $\Supp(M)=\Spec(B)$ (\sref[I]{I.9.1.13}).
Condition (ii) implies, by (14.3.2), that every irreducible component of $X$ containing $x$ dominates $Y$; in other words, the inverse image in $A$ of every minimal prime ideal $\mathfrak{p}_i$ of $B$ is $0$, which means that $t\notin\mathfrak{p}_i$ for all $i$.
Finally, (iii) means that $M/tM$ is a $(B/tB)$-module \emph{reduced} (3.2.2).
It suffices to show that under these conditions, $M$ is an $A$-module \emph{without torsion} ($\mathbf{0}_\mathbf{I}$,~6.3.4), and for this it evidently suffices to show that $t$ is an $M$-regular element; but this follows from (3.4.6.1).
\end{proof}

\begin{corollary}[15.2.3]
\label{IV.15.2.3}
Let $Y$ be a locally Noetherian prescheme, $f:X\to Y$ a morphism locally of finite type, $x$ a point of $X$, $y=f(x)$.
Suppose the following conditions are satisfied:
\begin{enumerate}[label=(\roman*)]
  \item $f$ is universally open at the generic points of the irreducible components of $f^{-1}(y)$ containing $x$.
  \item $f^{-1}(y)$ is geometrically reduced (on $k(y)$) at the point $x$ (4.6.9).
  \item The ring $\sh{O}_y$ is reduced.
\end{enumerate}
Under these conditions, $f$ is flat at the point $x$.
\end{corollary}

\begin{proof}
It suffices to apply \sref{IV.15.2.2} with $\sh{F}=\sh{O}_X$ (4.6.22).
\end{proof}

\begin{remarks}[15.2.4]
\label{IV.15.2.4}
\begin{enumerate}[label=(\roman*)]
  \item The first three conditions of \sref{IV.15.2.2} do not change when one replaces $f$ by $f_\mathrm{red}$; but if $\mathfrak{N}$ is the nilradical of a local Noetherian ring $A$, $A/\mathfrak{N}$ is flat over $A/\mathfrak{N}$ but not over $A$ itself when $\mathfrak{N}\neq 0$ (Bourbaki, \emph{Alg.~comm.}, chap.~II, \textsection~3, n\textsuperscript{o}~2, cor.~2 of prop.~5).
  One thus sees that one cannot suppress condition~(iv) in \sref{IV.15.2.2}.
  \item It follows from (2.4.6) that if the conclusion of \sref{IV.15.2.2} holds, then $f$ is universally open in a neighbourhood of $x$, thus in particular universally open at the points of the generic points of the irreducible components of $f^{-1}(y)$ containing $x$.
  Even when (i), (iii), and (iv) are satisfied, one cannot replace (ii) by the weaker hypothesis that $f$ is universally open at the point $x$ \emph{only}.
  This is what the example (14.1.3,~(i)) shows, where $f$ is evidently universally open at the point $x$, intersection of $X_1$ and $X_2$; the conditions (ii) and (iii) of \sref{IV.15.2.3} are also satisfied by this example.
  However, $\sh{O}_x$ is not an $\sh{O}_y$-module without torsion, $\sh{O}_y$ identifying with a subring of $\sh{O}_x$ that contains the divisors of zero; hence $f$ is not flat at the point $x$.
  \item In \sref{IV.15.2.3}, the conclusion is no longer necessarily valid when one replaces hypothesis~(ii) by the weaker hypothesis that $f^{-1}(y)$, considered as a $k(y)$-prescheme, \emph{has no associated embedded prime cycles}.
  An example is furnished by taking for $Y$ a curve having a cusp at a point, say $Y=\Spec(K[S,T]/(S^3-T^2))$, $K$ an algebraically closed field, and for $X$ its normalisation; the morphism $f:X\to Y$ is finite, surjective, and radiciel, hence a universal homeomorphism (2.4.5), but $f^{-1}(y)$ is not a reduced $k(y)$-prescheme.
  Here $f$ is a finite morphism, surjective, hence $f:X\to Y$ is universally open and a radiciel morphism.
  But $f$ is not flat at the point $x$, for if $\sh{O}_x$ were an $\sh{O}_y$-module flat, it would be an $\sh{O}_y$-module free of rank~$1$, which is absurd.
  \item
\oldpage[IV-3]{228}
  Let us show finally that in \sref{IV.15.2.2} or \sref{IV.15.2.3}, one cannot replace ``geometrically reduced'' by ``reduced''.
  We shall define two rings $A$, $A'$ having the following properties:
  \begin{enumerate}[label=\arabic{enumi})]
    \item $A'$ is the integral closure of $A$, an $A$-algebra finite whose residue field $k'$ is a non-trivial finite radiciel extension of the residue field $k=A/\mathfrak{m}$ of $A$.
    \item One can take $Y=\Spec(A)$, $X=\Spec(A')$, $\sh{F}=\sh{O}_x$, $y$ and $x$ being the closed points of $Y$ and $X$ respectively; the hypotheses (i) and (iv) of \sref{IV.15.2.2} are trivially satisfied; since $A'$ is a discrete valuation ring, $f:X\to Y$ is radiciel, finite, and surjective, hence a universal homeomorphism (2.4.5), which proves hypothesis~(i) of \sref{IV.15.2.3}.
  Finally it is clear that $f^{-1}(y)$ is reduced, but $f$ is not flat at the point $x$, for if $A'$ were an $A$-module flat, it would be an $A$-module free of rank~$1$ (Bourbaki, \emph{Alg.~comm.}, chap.~II, \textsection~3, n\textsuperscript{o}~2, prop.~5), which is absurd.
  \end{enumerate}
  \item The proof of \sref{IV.15.2.2} shows that one can weaken condition~(iii) by replacing ``geometrically reduced'' by ``reduced'' provided one only uses discrete valuation rings $A'$ whose residue field $k'$ is \emph{separable} over $k(y)$; in effect, if $x'$ and $z'$ have the same meanings as in \sref{IV.15.2.2}, the radiciel multiplicities $\mu_s(k(z')|k')$ and $\mu_s(k(z)|k(y))$ are the same, by (4.7.3), hence $\operatorname{long}((\sh{F}_{y'})_z)$ and $\operatorname{long}((\sh{F}_{y'})_{z'})$ are equal, and one can get the result still in the case $Y$ is the spectrum of a discrete valuation ring $y$ and $y$ its closed point.
  For example, one can obtain this result still when $\sh{O}_y$ is \emph{unibranch} and dominated by a discrete valuation ring whose residue field is a \emph{separable} extension of $k(y)$, by (11.6.4).
  This is the case when $\sh{O}_y$ is a \emph{regular} ring.
  But as Hironaka has pointed out, there are local rings of algebraic schemes, of dimension~$2$ over imperfect fields, that do not satisfy this condition.
\end{enumerate}
\end{remarks}

\subsection{Applications: criteria for reduction and irreducibility}
\label{subsection:IV.15.3}

\begin{proposition}[15.3.1]
\label{IV.15.3.1}
\oldpage[IV-3]{228}
Let $Y$ be a locally Noetherian prescheme, $f:X\to Y$ a morphism locally of finite type, $\sh{F}$ a coherent $\sh{O}_X$-Module, $x$ a point of $X$, $y=f(x)$.
Suppose that the conditions (i), (ii), and (iii) of \sref{IV.15.2.2} are satisfied.
Then:
\begin{enumerate}[label=(\roman*)]
  \item There exists a neighbourhood $U$ of $x$ in $X$ such that $f|U$ is universally open and, for every $x'\in U$, $\sh{F}_{f(x')}$ is geometrically reduced on $k(f(x'))$ at the point $x'$.
  \item If moreover $Y$ is reduced at the point $y$, there exists a neighbourhood $U_1\subset U$ of $x$ such that $\sh{F}|U_1$ is $f$-flat and reduced.
\end{enumerate}
\end{proposition}

\begin{proof}
By (\sref[I]{I.5.1.8}), (4.2.7), and (4.7.11), the hypotheses do not change, nor the conclusion~(i) of the statement, if one replaces $Y$ by $Y_\mathrm{red}$ and $X$ by $X\times_Y Y_\mathrm{red}$; one may thus suppose $Y$ \emph{reduced}.
It then follows from \sref{IV.15.2.2} that $\sh{F}$ is $f$-flat at the point $x$, hence also (11.1.1) in a neighbourhood of $x$, and \emph{a fortiori} $f$ is universally open in this neighbourhood (2.4.6).
The fact that $\sh{F}_{f(x')}$ is geometrically reduced at every point $x'$ of a neighbourhood of $x$ follows from (12.1.1,~(vii)).
Assertion~(ii) follows from (3.3.4).
\end{proof}

\begin{proposition}[15.3.2]
\label{IV.15.3.2}
The notation being that of \sref{IV.15.3.1}, suppose that the conditions (i), (ii), (iii) of \sref{IV.15.2.2} are satisfied, and suppose moreover that $f^{-1}(y)$ is equidimensional at the point $x$.
Then $f$ is equidimensional at the point $x$.
\end{proposition}

\begin{proof}
For every $y'\in Y$, $\Supp(\sh{F}_{y'})=f^{-1}(y')$.
By hypothesis, $\sh{F}_{y'}$ is reduced (hence satisfies condition~($\mathrm{S}_1$)) and is equidimensional at $x$.
By (12.1.1,~(ii)), we may suppose $\sh{F}_{f(x')}$ is equidimensional and of constant dimension $e=\dim_x(f^{-1}(y))$ at the points $x'$ for all $x'\in U$, which means that $f^{-1}(f(x'))$ is equidimensional at $x'$.
Since moreover $\sh{F}|U$ is $f$-flat, it follows from (2.3.4) and (13.3.1,~$a$) that $f$ is equidimensional at $x$.
\end{proof}

\begin{corollary}[15.3.3]
\label{IV.15.3.3}
The notation being that of \sref{IV.15.3.1}, suppose the following conditions satisfied:
\begin{enumerate}[label=(\roman*)]
  \item $\Supp(\sh{F})=X$.
  \item $f$ is universally open at the generic points of the irreducible components of $f^{-1}(y)$ containing $x$.
  \item $\sh{F}_y$ is geometrically punctually integral on $k(y)$ at the point $x$ (4.6.22).
  \item $Y$ is geometrically unibranch at the point $y$.
\end{enumerate}
Then $x$ belongs to a single irreducible component of $X$.

If moreover $Y$ is integral at the point $y$ and $\sh{F}=\sh{O}_X$, then $X$ is integral at the point $x$ and $f$ is flat at the point $x$.
\end{corollary}

\begin{proof}
Note first that (i) and (iii) imply that $x$ belongs to only a single irreducible component $Z$ of $f^{-1}(y)$.
It follows then from (14.4.7) and \sref{IV.15.3.2} that for every irreducible component $X_i$ of $X$ containing $x$ (hence $Z$), the restriction $f_i$ of $f$ to $X_i$ is equidimensional at the point $x$, and universally open at the generic point of $Z$.
The first assertion of the statement then follows from \sref{IV.15.1.3} and the second from \sref{IV.15.3.1}.
\end{proof}

\begin{remarks}[15.3.4]
\label{IV.15.3.4}
\begin{enumerate}[label=(\roman*)]
  \item The example (14.4.10,~(ii)) shows that in \sref{IV.15.3.3}, one cannot suppress the hypothesis that $Y$ is geometrically unibranch at the point $y$.
  \item The remark (15.2.4,~(vi)) shows that the conclusion of \sref{IV.15.3.3} is still valid under the following hypotheses: $\sh{O}_y$ is regular, $\Supp(\sh{F})=X$, and $\sh{F}_y$ is \emph{integral} at the point $x$.
  We do not know whether in this statement one can replace the hypothesis that $\sh{O}_y$ is regular by that of being geometrically unibranch, or even integral and integrally closed.
\end{enumerate}
\end{remarks}

\subsection{Complements on Cohen--Macaulay morphisms}
\label{subsection:IV.15.4}

\begin{proposition}[15.4.1]
\label{IV.15.4.1}
\oldpage[IV-3]{229}
Let $Y$ be a locally Noetherian prescheme, $f:X\to Y$ a morphism locally of finite type, $\sh{F}$ a coherent $\sh{O}_X$-Module such that $\Supp(\sh{F})=X$, $x$ a point of $X$, $y=f(x)$.
For every $z\in X$, set $X_{f(z)}=f^{-1}(f(z))$, $\sh{F}_{f(z)}=\sh{F}\otimes_{\sh{O}_{f(z)}}k(f(z))$.
Suppose that $\sh{F}$ is $f$-flat at the point $x$ and that $\sh{F}_y$ is an $\sh{O}_{X_y}$-Module de Cohen--Macaulay at the point $x$.
Then $f$ is equidimensional at the point $x$.
\end{proposition}

\begin{proof}
Indeed ((11.1.1) and (12.1.1,~(vi))), there exists a neighbourhood $U$ of $x$ such that $\sh{F}|U$ is $f$-flat and that for every $x'\in U$, $\sh{F}_{f(x')}$ is an $\sh{O}_{X_{f(x')}}$-Module de Cohen--Macaulay at the point $x'$; this last property implies that $\sh{F}_{f(x')}$ is equidimensional at $x'$, and that $x'$ does not belong to any embedded associated prime cycle of $\sh{F}_{f(x')}$ ($\mathbf{0}$,~16.5.4).
Moreover, by (12.1.1,~(ii)), one can suppose that the dimensions of the irreducible components of $\Supp(\sh{F}_{f(x')})|(U\cap X_{f(x')})$ take a (same) value independent of $x'$.
As $\Supp(\sh{F}_{f(x')})=X_{f(x')}$ by hypothesis, it follows from (2.3.4) and (13.3.1,~$a$) that $f$ is equidimensional at the point $x$.
\end{proof}

\begin{proposition}[15.4.2]
\label{IV.15.4.2}
Let $Y$ be a locally Noetherian prescheme, $f:X\to Y$ a morphism locally of finite type, $\sh{F}$ a coherent $\sh{O}_X$-Module such that $\Supp(\sh{F})=X$, $x$ a point of $X$, $y=f(x)$.
Suppose that $\sh{O}_y$ is regular and that $\sh{F}$ is an $\sh{O}_X$-Module de Cohen--Macaulay at the point $x$.
Then (with the notation of \sref{IV.15.4.1}) the following conditions are equivalent:
\begin{enumerate}[label=\alph*)]
  \item $\sh{F}$ is $f$-flat at the point $x$ and $\sh{F}_y$ is an $\sh{O}_{X_y}$-Module de Cohen--Macaulay at the point $x$.
  \item $\sh{F}$ is $f$-flat at the point $x$.
  \item $f$ is universally open in a neighbourhood of $x$.
  \item $f$ is open at the generic points of the irreducible components of $X_y$ containing $x$.
  \item $\dim(\sh{F}_x)=\dim(\sh{O}_y)+\dim((\sh{F}_y)_x)$.
  \item[e')] $\dim(\sh{O}_{X,x})=\dim(\sh{O}_y)+\dim(\sh{O}_{X,x}\otimes_{\sh{O}_y}k(y))$.
\end{enumerate}
\end{proposition}

\begin{proof}
As $\Supp(\sh{F})=X$, the conditions $e)$ and $e')$ are equivalent by definition (5.1.12).
Condition $a)$ implies trivially $b)$; $b)$ implies that $\sh{F}$ is $f$-flat at neighbouring points of $x$ (11.1.1), hence $b)$ implies $c)$ (2.4.6); $c)$ implies trivially $d)$; and $d)$ implies $e')$ by (14.2.1).
Finally, since $\sh{O}_y$ is regular and $\sh{F}_x$ is Cohen--Macaulay, it follows from (6.1.4) that $e)$ implies $a)$.
\end{proof}

\begin{proposition}[15.4.3]
\label{IV.15.4.3}
Let $f:X\to Y$ be a morphism locally of finite presentation, $\sh{F}$ an $\sh{O}_X$-Module of finite presentation such that $\Supp(\sh{F})=X$, $x$ a point of $X$, $y=f(x)$.
Suppose (with the notation of \sref{IV.15.4.1}) that the following conditions are satisfied:
\begin{enumerate}[label=\alph*)]
  \item $\sh{F}$ is $f$-flat at the point $x$ and $\sh{F}_y$ is an $\sh{O}_{X_y}$-Module de Cohen--Macaulay at the point $x$.
\end{enumerate}
Then:
\begin{enumerate}[label=\alph*)]
  \item There exists a neighbourhood $U$ of $x$ in $X$, and a $Y$-morphism quasi-finite $g:U\to Y'=Y[T_1,\ldots,T_n]$ ($T_i$ indeterminates), such that $\sh{F}|U$ is $g$-flat at $x$.
\end{enumerate}
\end{proposition}

\begin{proof}
Let us show first that $b)$ implies $a)$.
Set $h:Y'\to Y$; let $Y_y'=h^{-1}(y)$, the $k(y)$-prescheme regular ($\mathbf{0}$,~17.3.7); let $A$ be the local ring of $Y_y'$ at the point $y'=g(x)$, $k$ its residue field, and $B$ the local ring of $X_y$ at the point $x$; set moreover $(\sh{F}_y)_x=M$, which is a $B$-module of finite type.
The hypothesis that the morphism $g$ is quasi-finite implies in the same way that $g_y:X_y\to Y_y'$ (\textbf{II},~6.2.4), hence $M\otimes_A k$ is a $(B\otimes_A k)$-module of finite length (\textbf{II},~6.2.2); since by hypothesis $M$ is an $A$-module flat and $A$ is a Cohen--Macaulay ring (6.3.3), the fact that $\sh{F}$ is $f$-flat at $x$ results from that it is $g$-flat and the morphism $h$ is flat (2.1.6).

Let us now show that $a)$ implies $b)$.
The question being local on $X$ and $Y$, using (8.9.1) and (11.2.7), one reduces to the case where $X$ and $Y$ are Noetherian, using (6.7.1).
The hypothesis implies that $f$ is equidimensional at the point $x$ (\sref{IV.15.4.1}), whence the existence of a quasi-finite morphism $g:U\to Y'$ such that $g_y:U_y\to Y_y'$ is finite and surjective (13.3.1,~$b$); it suffices to prove that $\sh{F}|U$ is $g$-flat at $x$.
But by hypothesis $\sh{F}_y$ is $g_y$-flat at the point $x$.
But as $g_y$ is finite and surjective, it satisfies condition~$e')$ of \sref{IV.15.4.2} by (5.6.10); one may thus conclude by applying \sref{IV.15.4.2} to $g_y$ and $\sh{F}_y$.
\end{proof}


\subsection{Separable rank of fibres of a separated quasi-finite morphism and universally open morphism. Application to geometric connected components of fibres of a proper morphism.}
\label{subsection:IV.15.5}

\begin{proposition}[15.5.1]
\label{IV.15.5.1}
\oldpage[IV-3]{231}
Let $Y$ be a locally Noetherian prescheme, $f:X\to Y$ a separated and quasi-finite morphism.
For every $z\in Y$, let $n(z)$ be the number of geometric points of $f^{-1}(z)$ (or the \emph{separable rank} of $f^{-1}(z)$ on $k(z)$; cf.~\sref[I]{I.6.4.8}).
One has the following properties:
\begin{enumerate}[label=(\roman*)]
  \item The function $z\mapsto n(z)$ is semi-continuous inferiorly at every point $y\in Y$ such that $f$ is universally open at the points of $f^{-1}(y)$.
  \item Let $y$ be a point of $Y$ such that $f$ is universally open at the points of $f^{-1}(y)$; if the function $z\mapsto n(z)$ is constant in a neighbourhood of $y$, then $f$ is proper at the point $y$ (15.7.1).
  \item Suppose that $f$ is universally open and that every point of $Y$ is geometrically unibranch; then, if the function $z\mapsto n(z)$ is locally constant in $Y$, the irreducible components of $X$ are pairwise disjoint.
\end{enumerate}
\end{proposition}

\begin{corollary}[15.5.2]
\label{IV.15.5.2}
Under the general hypotheses of \sref{IV.15.5.1}, let $y$ be a point of $Y$ such that $f$ is universally open at the points of $f^{-1}(y)$.
Then, if $y'$ is a generisation of $y$, one has
\[
\tag{15.5.2.1}
  n(y) \leqslant n(y').
\]
\end{corollary}

\begin{lemma}[15.5.2.2]
\label{IV.15.5.2.2}
Under the general hypotheses of \sref{IV.15.5.1}, if $y$ is a point of $Y$, $y'$ a generisation of $y$ distinct from $y$, then there exists the spectrum of a discrete valuation ring $Z$, of closed point $z$ and generic point $z'$, and a morphism $g:Z\to Y$ such that: $1\textsuperscript{o}$ $g(z)=y$, $g(z')=y'$; $2\textsuperscript{o}$ if one sets $f'=f_{(Z)}$, $n(y)$ is the number of points of $f'^{-1}(z)$ and $n(y')$ is the number of points of $f'^{-1}(z')$.
\end{lemma}

\begin{proposition}[15.5.3]
\label{IV.15.5.3}
Let $Y$ be an irreducible locally Noetherian prescheme, $f:X\to Y$ a proper morphism.
For every $z\in Y$, let $n(z)$ be the number of geometric connected components of $f^{-1}(z)$.
Suppose that the restriction of $f$ to every irreducible component of $X$ is a dominant morphism.
Then, if $y_0$ is a point that is geometrically unibranch in $Y$, the function $y\mapsto n(y)$ is semi-continuous inferiorly at the point $y_0$.
\end{proposition}

\begin{proof}
Consider the Stein factorisation $f=g\circ f'$ of $f$ (\textbf{III},~4.3.3), where $g:Y'\to Y$ is a finite morphism and $f':X\to Y'$ a surjective morphism whose fibres are geometrically connected (\textbf{III},~4.3.4).
Since $f'$ is surjective, the inverse image by $f'$ of every irreducible component of $X$ contains at least an irreducible component of $Y'$; hence each of the irreducible components of $Y'$ dominates $Y$.
One may apply the Chevalley criterion (14.4.4) to the restrictions of $g$ to each of these irreducible components, and one concludes that $g$ is \emph{universally open at every point} of $g^{-1}(y_0)$.
If then $y$ is any point, since the $X_i\cap f^{-1}(y)$ are pairwise disjoint, $\sum_i n_i(y')=n(y')$; but by hypothesis the extremal members are equal, hence $n(y)=\sum_i n_i(y)$, which implies that the $X_i\cap f^{-1}(y)$ are pairwise disjoint.
The conclusion then follows from \sref{IV.15.5.2} and from the fact that the number of geometric connected components of $f^{-1}(y)$ equals the number of geometric points of $g^{-1}(y)$ (\textbf{III},~4.3.4).
\end{proof}

\begin{corollary}[15.5.4]
\label{IV.15.5.4}
Let $Y$ be a locally Noetherian prescheme, $f:X\to Y$ a proper morphism; let $y$ be a point of $Y$ such that $f$ is universally open at the points of $f^{-1}(y)$; then, with the notation of \sref{IV.15.5.3}, the function $z\mapsto n(z)$ is semi-continuous inferiorly at the point $y$.
\end{corollary}

\begin{remarks}[15.5.5]
\label{IV.15.5.5}
\begin{enumerate}[label=(\roman*)]
  \item Even if $Y$ is integral and normal, $X$ integral, $f$ finite and surjective (hence universally open by the Chevalley criterion (14.4.4)), the function $y\mapsto n(y)$ is not necessarily locally constant.
  \item The example (14.4.10,~(ii)) shows that in \sref{IV.15.5.1}, (iii), one cannot suppress the hypothesis that the points of $Y$ are geometrically unibranch.
  \item The example (11.7.5) shows that the conclusion of \sref{IV.15.5.1},~(i) is no longer valid if one only supposes the morphism $f$ is \emph{open} (even if $f$ is finite and surjective).
  \item Finally, in \sref{IV.15.5.1}, one cannot dispense with the hypothesis that the morphism $f$ is \emph{separated}, as the example of the affine line over a field with the origin ``doubled'' (\sref[I]{I.5.5.11}) shows: here $f$ is a local isomorphism (hence flat) and $Y$ is the affine line, but $z\mapsto n(z)$ is not semi-continuous inferiorly.
\end{enumerate}
\end{remarks}

\begin{lemma}[15.5.6]
\label{IV.15.5.6}
Let $Y$ be the spectrum of a discrete valuation ring, $a$ its closed point, $b$ its generic point.
Let $f:X\to Y$ be a morphism locally of finite type and open.
Suppose that $X$ is connected and that the fibre $f^{-1}(a)$ is a reduced prescheme.
Then $f^{-1}(b)$ is connected.
\end{lemma}

\begin{proof}
To prove \sref{IV.15.5.6}, we apply the lemma (15.5.6.1) to $X$ and to its closed part $Z=f^{-1}(a)$ which is rare since $f$ is open.
Note that by hypothesis, taking into account (15.2.2.1), $f$ is \emph{flat} at the points of $f^{-1}(a)$.
For such a point $x$, $\sh{O}_x$ is an $\sh{O}_a$-module without torsion ($\mathbf{0}_\mathbf{I}$,~6.3.4), and in particular, if $t$ is a uniformiser of $\sh{O}_a$, $t$ is $\sh{O}_x$-regular.
Moreover $\sh{O}_x/t\sh{O}_x$ is reduced by hypothesis; if it is of depth~$0$, it is therefore a field ($\mathbf{0}$,~16.4.7), and $t\sh{O}_x$ is then the maximal ideal of $\sh{O}_x$, hence $\Spec(\sh{O}_x)-\{x\}$ is reduced to a point, and \emph{a fortiori} connected.
If on the contrary $\operatorname{prof}(\sh{O}_x/t\sh{O}_x)\geqslant 1$, then $\operatorname{prof}(\sh{O}_x)\geqslant 2$, hence by the theorem of Hartshorne (5.10.7) $\Spec(\sh{O}_x)-\{x\}$ is connected, which finishes the proof.
\end{proof}

\begin{lemma}[15.5.6.1]
\label{IV.15.5.6.1}
Let $X$ be a locally Noetherian topological space, irreducible.
Let $Z$ be a closed rare part of $X$ having the following property: for every $x\in Z$, if one denotes by $\mathrm{T}_x$ the set of generisations of $x$ in $X$, then $\mathrm{T}_x-\{x\}$ is connected.
Under these conditions, if $X$ is connected, $X-Z$ is connected.
\end{lemma}

\begin{proposition}[15.5.7]
\label{IV.15.5.7}
\oldpage[IV-3]{235}
Let $Y$ be a locally Noetherian prescheme, $f:X\to Y$ a proper morphism.
For every $z\in Y$, let $n(z)$ be the number of geometric connected components of $f^{-1}(z)$.
Let $y$ be a point of $Y$ such that $f$ is universally open at the points of $f^{-1}(y)$ and that $f^{-1}(y)$ is geometrically reduced on $k(y)$ (4.6.2).
Then the function $z\mapsto n(z)$ is constant in a neighbourhood of $y$.
\end{proposition}

\begin{proof}
We already know that under the conditions of the statement, $z\mapsto n(z)$ is semi-continuous inferiorly superiorly, and by the same reasoning as at the beginning of the proof of \sref{IV.15.5.2}, it suffices to show that if $y'$ is a generisation of $y$, then
\[
\tag{15.5.7.1}
  n(y') \leqslant n(y).
\]
Using the reasoning from \sref{IV.15.5.2} and the lemma (15.5.2.2) applied to the Stein factorisation of $f$ (recalling that the number of geometric connected components of $f^{-1}(y)$ equals the number of geometric points of $g^{-1}(y)$ (\textbf{III},~4.3.4)), one reduces to the case where $n(y)$ and $n(y')$ are the numbers of points of $f^{-1}(y)$ and $f^{-1}(y')$ respectively, and where $Y$ is the spectrum of a discrete valuation ring, $y$ its closed point and $y'$ its generic point.
We may moreover replace $X$ by any connected component of its connected components, and these being both open and closed in $X$, proper on $Y$; the hypothesis $f$ is open at the points of $f^{-1}(y)$, the lemma \sref{IV.15.5.6} implies that $n(y')=1$, which completes the proof of (15.5.7.1).
\end{proof}

\begin{remark}[15.5.8]
\label{IV.15.5.8}
The relation (15.5.7.1) is no longer necessarily valid if $f$ is not supposed proper, even if the other hypotheses of \sref{IV.15.5.7} are satisfied.
With the notation of (15.5.5,~(i)), it suffices to consider the restriction of $f$ to $X-\{x'\}$, where $x'$ is one of the two points of $X$ above a point $y$ distinct from $y'$ and $y''$; then $n(y)=1$, while if $\eta$ is the generic point of $Y$, $n(\eta)=2$.
\end{remark}

\begin{proposition}[15.5.9]
\label{IV.15.5.9}
Let $f:X\to Y$ be a morphism flat and locally of finite presentation; for every $y\in Y$, let $n(y)$ be the number of geometric connected components of $f^{-1}(y)$.
\begin{enumerate}[label=(\roman*)]
  \item If $f$ is separated and quasi-finite, the function $y\mapsto n(y)$ (equal then to the number of geometric points of $f^{-1}(y)$) is semi-continuous inferiorly in $Y$; if it is constant in a neighbourhood of $y_0$, $f$ is proper at the point $y_0$.
  \item If $f$ is proper, the function $y\mapsto n(y)$ is semi-continuous inferiorly in $Y$; if moreover $f^{-1}(y_0)$ is geometrically reduced on $k(y_0)$, $n(y)$ is constant in a neighbourhood of $y_0$.
\end{enumerate}
\end{proposition}

\begin{proof}
In all cases, the questions are local on $Y$, so one may suppose $Y$ affine and $f$ of finite presentation; using (8.9.1), (8.10.5), and (11.2.7), one reduces to the case where $Y$ is Noetherian (using also (8.2.11) for the second assertion of~(i)).
Since moreover $f$ is then universally open (2.4.6), it suffices to apply \sref{IV.15.5.1}, (15.5.4), and \sref{IV.15.5.7} to conclude.
\end{proof}

\begin{remark}[15.5.10]
\label{IV.15.5.10}
One will note that another proof has thus been obtained of (12.2.4,~(vi)).
Conversely, one can deduce \sref{IV.15.5.7} from (12.2.4,~(vi)): in effect, on can restrict to the case where $Y$ is reduced and $f$ flat, and it follows then from the hypotheses of \sref{IV.15.5.7} and from \sref{IV.15.2.3} that $f$ is flat in a neighbourhood open of $f^{-1}(y)$; as moreover $f$ is proper, one can suppose this neighbourhood is of the form $f^{-1}(U)$, where $U$ is an open neighbourhood of $y$ in $Y$.
The proof of \sref{IV.15.5.7} given above has the advantage of putting in evidence the result (15.5.6), which is of independent interest.
\end{remark}


\subsection{Connected components of fibres along a section}
\label{subsection:IV.15.6}

\begin{proposition}[15.6.1]
\label{IV.15.6.1}
\oldpage[IV-3]{236}
Let $Y$ be a locally Noetherian prescheme, $f:X\to Y$ a morphism locally of finite type, $g:Y\to X$ a $Y$-section of $X$ (\sref[I]{I.2.5.5}).
For every $y\in Y$, denote by $X_y^0$ the connected component of $f^{-1}(y)$ containing the point $g(y)$.
Let $y$ be a point of $Y$, $v'$ a generisation of $y$.
Then:
\begin{enumerate}[label=(\roman*)]
  \item If there exists an irreducible component $Z$ of $X_y^0$, of dimension equal to $\dim(X_y^0)$ and containing $g(y')$ (which will be the case if $X_y^0$ is irreducible), then
\[
\tag{15.6.1.1}
  \dim(X_y^0) \geqslant \dim(X_{y'}^0).
\]
  \item If moreover $X_y^0$ is irreducible and the two members of (15.6.1.1) are equal, then $X_{y'}^0\subset\overline{X_y^0}$ (closure in $X$).
\end{enumerate}
\end{proposition}

\begin{proof}
(i) Let $\overline{Z}$ be the closure of $Z$ in $X$; since $y$ is adherent to $y'$ and $g$ continuous, $g(y)\in\overline{Z}$.
As $Z=\overline{Z}\cap f^{-1}(y')$, it follows from Chevalley's semi-continuity theorem (13.1.3) applied to the restriction of $f$ to the closed reduced sub-prescheme of $X$ having $\overline{Z}$ as underlying space, that $\dim_{g(y)}(\overline{Z}\cap f^{-1}(y))\geqslant\dim(Z)$; but the irreducible components of $\overline{Z}\cap f^{-1}(y)$ containing $g(y)$ are evidently contained in $X_y^0$, whence the inequality (15.6.1.1).

(ii) The reasoning of~(i) shows moreover that if the two members of (15.6.1.1) are equal, $\dim_{g(y)}(\overline{Z}\cap f^{-1}(y'))=\dim_{g(y)}(X_y^0)$; if $X_y^0$ is irreducible, then $\overline{Z}=X_y^0$, hence $X_y^0$ is contained in $\overline{Z}$, and \emph{a fortiori} in $\overline{X_y^0}$.
\end{proof}

\begin{corollary}[15.6.2]
\label{IV.15.6.2}
With the notation of \sref{IV.15.6.1}, suppose moreover that for every $y\in Y$, $X_y^0$ is irreducible.
Then the function $z\mapsto\dim(X_z^0)$ is semi-continuous superiorly in $Y$.
If moreover this function is constant in a neighbourhood of a point $y\in Y$, then $X_y^0\subset\overline{X_{y'}^0}$ for every generisation $y'$ of $y$.
\end{corollary}

\begin{proof}
Let $y$ be a point of $Y$, and let $U$ be an open affine neighbourhood of $g(y)$ in $X$; then for every $z\in V$, $U\cap X_z^0$ is dense in $X_z^0$, hence $\dim(X_z^0)=\dim(U\cap X_z^0)$ (4.1.3).
Replacing $X$ by $U$, one can thus suppose $f$ is of finite type.
One then knows (9.7.10) that the function $z\mapsto\dim(X_z^0)$ is locally constructible, and the assertions of the corollary follow from \sref{IV.15.6.1} and ($\mathbf{0_{III}}$,~9.3.4).
\end{proof}

\begin{proposition}[15.6.3]
\label{IV.15.6.3}
With the notation of \sref{IV.15.6.1}, suppose that for every $y\in Y$, $X_y^0$ is geometrically irreducible (4.5.2).
Then, for every $y\in Y$ such that the function $z\mapsto\dim(X_z^0)$ is constant in a neighbourhood of $y$, $f$ is universally open at $X_y^0$.
\end{proposition}

\begin{proof}
Apply the criterion (14.3.7).
Let $h:Y'\to Y$ be a morphism, $Y'$ being the spectrum of a discrete valuation ring, and denote by $t$, $t'$ the closed point and the generic point respectively, and suppose $h(t)=y$, and hence $y'=h(t')$ is a generisation of $y$.
Set $X'=X\times_Y Y'$, $f'=f_{(Y')}:X'\to Y'$, $g'=g_{(Y')}:Y'\to X'$.
With the same notation as in \sref{IV.15.6.1}, the hypothesis on the $X_y^0$ implies that $X_t'^0=X_y^0\otimes_{k(y)}k(t)$ and $X_{t'}'^0=X_{y'}^0\otimes_{k(y')}k(t')$ are irreducible (4.4.1); since $\dim(X_y^0)=\dim(X_{y'}^0)$ and $\dim(X_y^0)=\dim(X_t'^0)$ (4.1.4), we see that one is reduced to proving the proposition for $f'$ and $X_t'^0$, i.e.\ one can restrict to the case where $Y$ is the spectrum of a discrete valuation ring, $y$ its closed point and $y'$ its generic point; if $x'$ is the generic point of $X_y^0$, it follows then from \sref{IV.15.6.2} that the point of $X_{y'}^0$ is a specialisation of $x'$, whence the conclusion.
\end{proof}

\begin{proposition}[15.6.4]
\label{IV.15.6.4}
Under the general hypotheses of \sref{IV.15.6.1}, let $X^0$ be the union of the $X_y^0$ when $y$ ranges over $Y$.
Suppose that for every $y\in Y$, $X_y^0$ is geometrically reduced on $k(y)$ (4.6.2), and that $f$ is universally open at every point of $X_y^0$.
Then $X^0$ is a neighbourhood of $X_y^0$ in $X$.
In particular, if $f$ is universally open and if, for every $y\in Y$, $X_y^0$ is geometrically reduced on $k(y)$, then $X^0$ is open in $X$.
\end{proposition}

\begin{proof}
One can first reduce to the case where $f$ is a morphism of \emph{finite type}.
Suppose then that $f$ is of finite type.
Then $X^0$ is locally constructible (9.7.12); it suffices to prove that $X^0$ contains every generisation $x'$ of $x$ ($\mathbf{0_{III}}$,~9.2.5).
Set $y'=f(x')$; there exists a discrete valuation ring $A$ and a morphism $u$ of $Y'=\Spec(A)$ to $X$ such that if $z$ is the closed point and $z'$ the generic point, $u(z)=x$ and $u(z')=x'$ (\textbf{II},~7.1.9); if $h=f\circ u$, $h(z)=y$ and $h(z')=y'$.
Set $X'=X\times_Y Y'$, $f'=f_{(Y')}:X'\to Y'$, $g'=g_{(Y')}:Y'\to X'$.
If $p_y:f'^{-1}(z)\to f^{-1}(y)$ (resp.\ $p_{y'}:f'^{-1}(z')\to f^{-1}(y')$) is the canonical projection, $p_y^{-1}(X_y^0)$ (resp.\ $p_{y'}^{-1}(X_{y'}^0)$) is the connected component of $f'^{-1}(z)$ (resp.\ $f'^{-1}(z')$) containing $g'(z)$ (resp.\ $g'(z')$), since the $X_i^0$ are geometrically connected (4.5.13), and it suffices to apply (4.4.1).
Moreover, the morphism $u$ corresponds to a $Y'$-section $v$ of $X'$ such that $p_y(v(y))=x$, $p_y(v(y'))=x'$; as $v(y')$ is a generisation of $v(y)$, and $f'^{-1}(y)-X_y^0=f^{-1}(y)$ is then closed in $f^{-1}(y)$, we can, replacing $X$ by the open $X-(f^{-1}(y)-X_y^0)$, suppose $X^0=f^{-1}(y)$, i.e.\ $X^0=X$; in other words, the proposition will be established if one shows that $X^0=X$, or yet again that $X_{y'}$ is connected; but $f^{-1}(y')$ is \emph{reduced}; as moreover $f$ is open at the points of $f^{-1}(y)$, the lemma \sref{IV.15.5.6} allows us to conclude.
\end{proof}

\begin{corollary}[15.6.5]
\label{IV.15.6.5}
Let $f:X\to Y$ be a morphism flat and of finite presentation, $g$ a $Y$-section of $X$; for every $y\in Y$, let $X_y^0$ be the connected component of $f^{-1}(y)$ containing $g(y)$, and let $X^0$ be the union of the $X_y^0$ when $y$ ranges over $Y$.
Then, for every $y\in Y$ such that $X_y^0$ is geometrically reduced on $k(y)$, $X^0$ is a neighbourhood of $X_y^0$ in $X$.
In particular, if $f$ is reduced (6.8.1), $X^0$ is open in $X$.
\end{corollary}

\begin{proof}
The question being local on $Y$, one can suppose $Y=\Spec(A)$ affine.
There exists then a sub-Noetherian ring $A_0$ and a morphism of finite type $f_0:X_0\to Y_0=\Spec(A_0)$ such that $X=X_0\times_{Y_0}Y$ and $f=(f_0)_{(Y)}$ ((8.9.1) and (11.2.7)); moreover, one can suppose there exists a $Y_0$-section $g_0$ of $X_0$ such that $g=(g_0)_{(Y)}$ (8.8.2).
For every $y\in Y$, let $y_0$ be its projection in $Y_0$; it follows from (4.5.13) that $(X_0)_{y_0}^0$ is geometrically connected, hence if $p:f^{-1}(y)\to f_0^{-1}(y_0)$ is the canonical projection, $X_y^0=p^{-1}((X_0)_{y_0}^0)$ (4.4.1); moreover, if $f^{-1}(y)$ is geometrically reduced on $k(y)$, it is also geometrically reduced on $k(y_0)$ (4.6.10).
On thus reduces to the case where $Y$ is Noetherian and $f$ of finite type.
Since moreover $f$ is then universally open (2.4.6), (11.1.1), it suffices to apply \sref{IV.15.6.4}.
\end{proof}

\begin{proposition}[15.6.6]
\label{IV.15.6.6}
\oldpage[IV-3]{239}
Let $f:X\to Y$ be a morphism such that $Y$ is locally Noetherian and $f$ locally of type finite, or that $f$ is locally of finite presentation.
Let $g$ be a $Y$-section of $X$; for every $y\in Y$, let $X_y^0$ be the connected component of $f^{-1}(y)$ containing $g(y)$; finally, let $X^0$ be the union of the $X_y^0$ when $y$ ranges over $Y$.

Suppose that for every $y\in Y$, $X_y^0$ is geometrically irreducible, and that $f$ is universally open at $X_y^0$.
Then:
\begin{enumerate}[label=(\roman*)]
  \item The function $y\mapsto\dim(X_y^0)$ is semi-continuous superiorly in $Y$.
  \item Let $x_0\in X^0$, $y_0=f(x_0)$.
  If the function $y\mapsto\dim(X_y^0)$ is constant in a neighbourhood of $y_0$, then there exists a neighbourhood open $U$ of $y_0$ in $Y$ such that $f$ is universally open at the points of $X^0\cap f^{-1}(U)$, and $f$ is flat at the point $x_0$.
  \item Conversely, suppose that $f$ is universally open at $X_{y_0}^0$, and that moreover one of the following conditions is satisfied:
  \begin{enumerate}[label=$\alpha$)]
    \item The fibre $f^{-1}(y_0)$ is geometrically reduced on $k(y_0)$ at the point $g(y_0)$ and the ring $\sh{O}_{Y,y_0}$ is Noetherian.
    \item For the generic point $y'$ of an irreducible component of $Y$ containing $y_0$, one has $\dim(f^{-1}(y'))=\dim(X_{y'}^0)$.
  \end{enumerate}
  Then the function $y\mapsto\dim(X_y^0)$ is constant in a neighbourhood of $y_0$.
  \item The set $W$ of the $x\in X^0$ such that $y\mapsto\dim(X_y^0)$ is constant in a neighbourhood of $f(x)$ and that $X_{f(x)}^0$ is geometrically reduced on $k(f(x))$, is open in $X$.
\end{enumerate}
\end{proposition}

\begin{corollary}[15.6.7]
\label{IV.15.6.7}
Suppose the preliminary conditions of \sref{IV.15.6.6} on $f$ satisfied, and suppose moreover that for every $y\in Y$, $X_y^0$ is geometrically integral on $k(y)$.
Then the following conditions are equivalent:
\begin{enumerate}[label=\alph*)]
  \item The function $y\mapsto\dim(X_y^0)$ is locally constant in $Y$.
  \item The morphism $f_{(Y_\mathrm{red})}:X\times_Y Y_\mathrm{red}\to Y_\mathrm{red}$ deduced from $f$ by base change is flat at the points of $X^0$.
\end{enumerate}

Moreover, these conditions imply that $X^0$ is open in $X$.
Finally, when $Y$ is locally Noetherian, $a)$ and $b)$ are also equivalent to:
\begin{enumerate}
  \item[b')] $f$ is universally open at the points of $X^0$.
\end{enumerate}
\end{corollary}

\begin{proposition}[15.6.8]
\label{IV.15.6.8}
Let $f:X\to Y$ be a separated morphism of finite presentation, $g$ a $Y$-section of $X$; for every $y\in Y$, let $X_y^0$ be the connected component of $f^{-1}(y)$ containing $g(y)$, and let $X^0$ be the union of the $X_y^0$ when $y$ ranges over $Y$.
If for every $y\in Y$ $X_y^0$ is proper on $k(y)$ (i.e.~(15.7.1), there exists a neighbourhood open $U$ of $y$ in $Y$ such that $X^0\cap f^{-1}(U)$ is closed in $f^{-1}(U)$ and proper on $U$).

Proceeding as at the beginning of the proof of \sref{IV.15.6.6}, and using (8.10.5,~(v)), and replacing (9.7.11) by (9.7.7), one reduces to the case where $Y$ is Noetherian and $f$ of finite type.
One then knows (9.7.12), (9.7.11) that $X^0$ is locally constructible in $X$; on the other hand, the $X_y^0$ are geometrically connected (for every $z\in Y$) by (4.5.13); finally, $X^0$ is universally submersive over $Y$ (15.7.8).
It suffices then to apply the criterion (15.7.9).
\end{proposition}

\begin{corollary}[15.6.9]
\label{IV.15.6.9}
Let $f:X\to Y$ be a separated morphism such that $Y$ is locally Noetherian and $f$ locally of finite type, or that $f$ is of finite presentation.
Let $g$ be a $Y$-section of $X$, and for every $y\in Y$, let $X_y^0$ be the connected component of $f^{-1}(y)$ containing $g(y)$.
Let $y$ be a point of $Y$ such that $X_y^0$ is proper on $k(y)$.
Then, for every generisation $y'$ of $y$, one has $\dim(X_y^0)\geqslant\dim(X_{y'}^0)$.
\end{corollary}

\begin{proof}
Suppose first $f$ of finite presentation; then replacing $Y$ by a neighbourhood of $y$ if necessary, one can suppose that $f|X^0$ is \emph{proper} \sref{IV.15.6.8}, and it suffices to apply (13.1.5) to this restriction.
If $Y$ is locally Noetherian and $f$ locally of finite type, $X$ is locally Noetherian; moreover the hypothesis implies that $X_y^0$ is Noetherian; therefore there is an open Noetherian $V\subset X$ containing $X_y^0$; the restriction $f$ to $V\cap f^{-1}(W)$ is a morphism of type finite, and $g|W$ is a $W$-section of $X$, and by applying the above result to this morphism, one gets $\dim(X_y^0)\geqslant\dim(X_{y'}^0)$.
\end{proof}

\begin{remarks}[15.6.10]
\label{IV.15.6.10}
\begin{enumerate}[label=(\roman*)]
  \item The supplementary hypothesis on the existence of an irreducible component of $X_y^0$ containing $g(y')$ and of dimension equal to $\dim(X_y^0)$, made in \sref{IV.15.6.1}, is not superfluous, as the following example shows.
  Let $A$ be a discrete valuation ring, $\pi$ a uniformiser of $A$, $K$ the field of fractions of $A$, $k$ its residue field.
  Consider the two $Y$-schemes $X_1=\Spec(A[t])$, $X_2=\Spec(K[u,v])$, where $t$, $u$, $v$ are indeterminates.
  The projection of $X_1$ in $Y$ has for image $\{y\}$ (i.e.\ the projection of $X_2$ to $Y$ is the point $\{y'\}$).
  One can \emph{glue} the two sub-preschemes $X_1$ and $X_2$ along their common open $Y$-isomorphism $Z_1\cong Z_2$, getting a $Y$-scheme $X$; $f^{-1}(y)$ is the disjoint union of $(X_1)_y$ and $(X_2)_y$, whilst $f^{-1}(y')$ has two irreducible components, isomorphic respectively to $(X_1)_{y'}$ and $(X_2)_{y'}$, having a unique common point $x$, which projects to the generic point $y'$ of $Y$ and such that $k(x_1)=K$.
  The ``null section'' $g$ of $X_1$ (\textbf{II},~1.7.9) is then a $Y$-section of $X$; $f^{-1}(y)$ equals $(X_1)_y$, while $f^{-1}(y')$ has two irreducible components isomorphic to $(X_1)_{y'}$ and $(X_2)_{y'}$, with a single common point; even though $\dim(X_{y'}^0)=1$ whilst $X_y^0$ is not a neighbourhood of $X_{y_0}^0$ in $X$ (one cannot suppress the hypotheses $\alpha$) and $\beta$) in \sref{IV.15.6.6}).
  Moreover, $X^0$ is not open in $X$, which shows that in \sref{IV.15.6.4} one cannot dispense with the hypothesis that $X'$ is reduced.
  \item Consider the morphism $f$ defined in (12.2.3,~(iii)), which is proper and flat; moreover the reciprocal $g:Y\to X$ is a $Y$-section of $X$.
  It satisfies $X_1\cap f_0^{-1}(Y)$ an isomorphism, hence the reciprocal $g:Y\to X$ is a $Y$-section of $X$.
  We have $\dim(X_y)=1$ whilst $\dim(X^0)=0$ for $y\neq y_0$, even though $X_y^0$ is not a neighbourhood of $X_{y_0}^0$ in $X$.
  Moreover, $X^0$ is not open in $X$, which shows that in \sref{IV.15.6.4}, one cannot dispense with the hypothesis that $X'$ is reduced.
  \item At chap.~VI, we shall apply the preceding results to $Y$-preschemes that are \emph{groups} locally of finite type.
  If $G$ is such a prescheme, there is always a $Y$-section canonical $g$, the ``unit section'' and one shows that $G_y$ is always geometrically irreducible and that one has $\dim(G_y^0)=\dim(G_y)$, by setting $G_y^0=f^{-1}(y)$.
  It follows then from \sref{IV.15.6.2} and \sref{IV.15.6.3} that the function $z\mapsto\dim(G_z)$ is semi-continuous superiorly, and that if it is constant in a neighbourhood of a point $y$, $f$ is universally open at all the points of $G_y^0$.
  If moreover $G_y$ is geometrically reduced on $k(y)$ at all its points, one shows that $G$ is \emph{smooth} (6.8.1) at all the points of $G_y^0$; it will follow then from \sref{IV.15.6.6} that if moreover $\sh{O}_y$ is reduced, then $f$ is \emph{smooth} at all the points of $G_y^0$.
  These results apply for example in the theory of Picard schemes, where one will have at our disposal a general theorem assuring that, under certain conditions, the function $z\mapsto\dim(G_z)$ is locally constant.
  We also note that in the case of a $Y$-prescheme in groups $G$, the hypothesis $\beta$) of \sref{IV.15.6.6} is always satisfied.

  One will note also that in the case of a $Y$-group prescheme $G$, the hypothesis $\beta$) of \sref{IV.15.6.6} is always satisfied.
\end{enumerate}
\end{remarks}


\subsection{Appendix: Valuative criteria for local properness}
\label{subsection:IV.15.7}

\begin{definition}[15.7.1]
\label{IV.15.7.1}
\oldpage[IV-3]{242}
Let $f:X\to Y$ be a morphism of preschemes, $y$ a point of $Y$.
One says that $f$ is \emph{proper at the point $y$} if there exists a neighbourhood open $U$ of $y$ in $Y$ such that the restriction $f^{-1}(U)\to U$ of $f$ is a proper morphism.
One says that a part $Z$ of $X$ is \emph{proper over $Y$ at the point $y$} if there exists a neighbourhood open $U$ of $y$ such that $Z\cap f^{-1}(U)$ is closed in $f^{-1}(U)$ and proper on $U$ (\textbf{II},~5.4.10) (which amounts to saying that for every sub-prescheme having $Z$ as closed underlying space, the restriction $Z\to Y$ of $f$ to $Z$ is proper at the point $y$).

Let $g:Y'\to Y$ be any morphism; set $X'=X\times_Y Y'$, $f'=f_{(Y')}$, and, if $p:X'\to X$ is the canonical projection, $Z'=p^{-1}(Z)$.
Then if $Z$ is proper on $Y$ at the point $y$, $Z'$ is proper on $Y'$ at every point $y'$ above $y$ (\textbf{II},~5.4.2).
\end{definition}

\begin{proposition}[15.7.2]
\label{IV.15.7.2}
Let $Y$ be an integral locally Noetherian prescheme, $X$ an integral prescheme, $f:X\to Y$ a separated morphism, dominant, and of finite type, $y$ a point of $Y$.
The following conditions are equivalent:
\begin{enumerate}[label=\alph*)]
  \item $f$ is proper at the point $y$.
  \item If $Y_1=\Spec(\sh{O}_y)$, $X_1=X\times_Y Y_1$, the morphism $f_1=f_{(Y_1)}:X_1\to Y_1$ is proper.
  \item Every discrete valuation ring having $\mathrm{R}(X)$ for field of fractions and dominating a local ring $\sh{O}_x$ of $X$ (in which case one has necessarily $f(x)=y$).
\end{enumerate}
\end{proposition}

\begin{proof}
The fact that $a)$ implies $b)$ follows from (\textbf{II},~5.4.2), and the implication $b)\Rightarrow c)$ follows from (\textbf{II},~7.3.10).
It remains to show that $c)$ implies $a)$.
The question being local on $Y$, one may suppose $Y$ affine, hence Noetherian.
By the lemma of Chow (\textbf{II},~5.6.1), there exists an integral prescheme $X'$, a projective morphism $p:P\to Y$, an open immersion $j:X'\to P$ dominant, and a projective morphism birational (hence surjective) $g:X'\to X$ such that the diagram
\[
\xymatrix{
P & X' \ar@{_{(}->}[l]_-{j} \ar[d]^{g}\\
Y & X \ar[l]
}
\]
is commutative.
Since $X'$ is integral and $j$ dominant, $P$ is irreducible, and one may, replacing $p$ by $p_\mathrm{red}$, suppose $P$ integral.
It suffices to prove that there exists an open $U$ of $y$ in $Y$ such that $f(X)\cap U$ is closed in $U$, and proper on $U$.
But the image $f(X)$ is locally closed in $Y$, since $f$ is of finite type, and open in its closure.
Since $f$ is of finite type, it factorises as $f=j_0\circ g_0$, where $j_0:Z'\to Y$ is the canonical injection and $g_0$ is the restriction $g_0(X)=Z''$ (which is open from $Z'$); the fact that $g$ is proper implies $d)$.
Finally, by the Chow lemma (\textbf{II},~7.2.2 and \textbf{II},~7.3.5,~(vi)), the uniqueness of the prolongation of a $Y'$-section follows from the hypothesis; by \sref{IV.15.7.2}, $X$ is proper over $Y$, which proves the proposition.
\end{proof}

\begin{remark}[15.7.3]
\label{IV.15.7.3}
One can, in \sref{IV.15.7.2}, suppress the hypothesis that $Y$ is locally Noetherian, in replacing in the condition $c)$ the discrete valuation rings by any valuation rings; the proof is then unchanged, using (\textbf{II},~7.3.10).
\end{remark}

\begin{corollary}[15.7.4]
\label{IV.15.7.4}
Let $Y$ be a locally Noetherian prescheme, $f:X\to Y$ a separated morphism, $Z$ a part of $X$ such that the maximal points of its closure $\overline{Z}$ belong to $Z$ (which is the case when $Z$ is \emph{finite}, or when $Z$ is \emph{constructible} ($\mathbf{0_{III}}$,~9.2.2)).
Let $y$ be a point of $Y$.
The following conditions are equivalent:
\begin{enumerate}[label=\alph*)]
  \item The set $\overline{Z}$ is proper over $Y$ at the point $y$.
  \item If $Y_1=\Spec(\sh{O}_y)$, $X_1=X\times_Y Y_1$, and if $g_1:X_1\to X$ is the canonical projection and $Z_1=g_1^{-1}(Z)$, the closure $\overline{Z}_1$ of $Z_1$ in $X_1$ is proper on $Y_1$.
  \item For every scheme $Y'=\Spec(A')$, where $A'$ is a discrete valuation ring, and every morphism $g:Y'\to Y$, such that the image $g(y')$ of the closed point $y'$ of $Y'$ is $y$, one has the following property: setting $X'=X\times_Y Y'$, $f'=f_{(Y')}:X'\to Y'$, $g'=g_{(X)}:X'\to X$, $Z'=g'^{-1}(Z)$, and denoting by $\eta'$ the generic point of $Y'$, then for every $x'\in Z'\cap f'^{-1}(\eta')$, rational on $k(\eta')$, there exists a $Y'$-section of $X'$ containing $x'$.
\end{enumerate}
\end{corollary}

\begin{proof}
The fact that $a)$ implies $b)$ follows from the fact that $\overline{Z}_1\subset g_1^{-1}(\overline{Z})$ and that $g_1^{-1}(Z)$ is proper on $Y_1$ \sref{IV.15.7.1}.
The fact that $b)$ implies $c)$ follows from (\textbf{II},~7.3.3).
It remains to show that $c)$ implies $a)$.
It suffices to prove that every irreducible component of $\overline{Z}$ is proper over $Y$ at the point $y$, and the hypothesis thus allows us to restrict to the case where $Z$ is reduced to a single point $z$.
One can also, by (\textbf{I},~5.2.2) and (\textbf{II},~5.4.6), suppose that $X$ and $Y$ are reduced, then (by definition of a proper part (\textbf{II},~5.4.10)) suppose that $\overline{Z}=X=\{\overline{z}\}$ and (applying again (\textbf{I},~5.2.2)) that $f$ is dominant, so $X$ and $Y$ are integral and Noetherian; one must then prove that $f$ is proper at the point $y$ of $Y$.
Let us show for this that $f$ satisfies condition $c)$ of \sref{IV.15.7.2}.
Let $A'$ be a discrete valuation ring having $\mathrm{R}(X)=k(z)$ for field of fractions and dominating $\sh{O}_y$; with the notations of $c)$, one has $g(y')=y$, and $g(\eta')$ is the generic point $\eta=f(z)$ of $Y$.
Since by definition $k(\eta')=k(z)$, there exists a $k(\eta)$-morphism $\Spec(k(\eta'))\to\Spec(k(z))$ sending $\eta'$ to $z$, hence a $k(\eta')$-section of $f'^{-1}(\eta')$ (\textbf{I},~3.3.14); if $z'$ is the image of $\eta'$ by this section, $z'$ is rational over $k(\eta')$ and $g'(z')=z$.
One concludes from hypothesis $c)$ that there exists a $Y'$-section $h'$ of $X'$ such that $h'(\eta')=z'$; let $h=g'\circ h':Y'\to X$ be the corresponding $Y$-morphism, and set $x=h(y')$; it is clear that $f(x)=y$ and that $A'=\sh{O}_{y'}$ dominates $\sh{O}_z$, which completes the proof.
\end{proof}

\begin{corollary}[15.7.5]
\label{IV.15.7.5}
Let $Y$ be a locally Noetherian prescheme, $f:X\to Y$ a separated morphism of finite type, $y$ a point of $Y$.
The following conditions are equivalent:
\begin{enumerate}[label=\alph*)]
  \item $f$ is proper at the point $y$.
  \item If $Y_1=\Spec(\sh{O}_y)$, $X_1=X\times_Y Y_1$, the morphism $f_1=f_{(Y_1)}:X_1\to Y_1$ is proper.
  \item For every scheme $Y'=\Spec(A')$, where $A'$ is a discrete valuation ring, and every morphism $g:Y'\to Y$, such that the image $g(y')$ of the closed point $y'$ of $Y'$ is $y$, one has the following property: setting $X'=X\times_Y Y'$, $f'=f_{(Y')}:X'\to Y'$, and denoting by $\eta'$ the generic point of $Y'$, for every $x'\in f'^{-1}(\eta')$, rational on $k(\eta')$, there exists a $Y'$-section of $X'$ containing $x'$.
\end{enumerate}
\end{corollary}

\begin{proof}
This is a special case of \sref{IV.15.7.4} applied to $Z=X$.
\end{proof}

\begin{corollary}[15.7.6]
\label{IV.15.7.6}
Let $Y$ be a locally Noetherian prescheme, $f:X\to Y$ a morphism of finite type.
The following conditions are equivalent:
\begin{enumerate}[label=\alph*)]
  \item There exists a proper morphism $g:X\to Z$ and an immersion $h:Z\to Y$ such that $f=h\circ g$.
  \item The sub-space $f(X)$ is locally closed in $Y$, and if $Z'$ is the sub-prescheme of $Y$ induced by the closure $\overline{f(X)}$ of $Z'$ on the open $f(X)$ of $Z'$, and $Z''$ the prescheme induced by $Z'$ on $f(X)$, then $f$ factorises uniquely as $f=j_0\circ g_0$, where $j_0:Z''\to Y$ is the canonical injection, and $g_0:X\to Z''$ is proper.
  \item For every $y\in f(X)$, $f$ is proper at the point $y$.
  \item For every scheme $Y'=\Spec(A')$, where $A'$ is a discrete valuation ring, and every morphism $g:Y'\to Y$, such that $g(Y')\subset f(X)$, there exists a $Y'$-section of $X'=X\times_Y Y'$, and every $Y'$-section rationale (\sref[I]{I.7.1.2}) of $X'\times_Y Y'$ extends uniquely to a $Y'$-section of $X'$.
\end{enumerate}
\end{corollary}

\begin{proof}
It is clear that $b)$ implies $a)$.
To see that $a)$ implies $c)$, note first that $a)$ implies that $h(Z)$ is locally closed in $Y$ and $g(X)$ is closed in $Z$, hence $f(X)$ is locally closed in $Y$; for every $y\in f(X)$, there is thus an open neighbourhood $U$ of $y$ in $Y$ such that $f(X)\cap U$ is closed in $U$, hence $f(X)$ is locally closed, and therefore open in its closure.
Since this closure is the underlying space of $Z'$, and $f$ factorises as $f=j_0\circ g_0$, where $j_0:Z'\to Y$ is the canonical injection and $g_0$ is a morphism from $X$ to $Z'$, the fact that $g_0(X)=Z''$ (which is open in $Z'$) implies the factorisation of the statement of $b)$; the fact that $g$ is proper then follows from the local character (on $Z''$) of this property, and from (\textbf{II},~5.4.3).
It is clear that $c)$ implies $d)$ by virtue of \sref{IV.15.7.5}, the uniqueness of the extension of a $Y'$-section resulting from the hypothesis that $f$ is separated (\textbf{I},~7.2.2).
Finally, let us prove that $d)$ implies $c)$; first, it follows from $d)$, by (\textbf{II},~7.2.3 and 7.2.4), that $f$ is separated; one can then apply the criterion (\sref{IV.15.7.5},~$c$)), which shows that $f$ is proper at every point of $f(X)$.
\end{proof}

\begin{remarks}[15.7.7]
\label{IV.15.7.7}
\begin{enumerate}[label=(\roman*)]
  \item In (15.7.4,~$c$) and (15.7.5,~$c$), one can restrict to the case where the discrete valuation ring $A'$ is \emph{complete} and has an algebraically closed residue field.
  Indeed, in the proof of \sref{IV.15.7.4}, if one considers a complete discrete valuation ring $A''$, dominant $A'$, and having an algebraically closed residue field, the hypothesis $c)$ is then satisfied for $Y''=\Spec(A'')$, $X''=X'\times_{Y'}Y''$, and $f''=f'_{(Y'')}$; but $X''=X'\times_{Y'}Y''$, and it follows then from the remark (\textbf{II},~7.3.9,~(i)) that $f'$ is proper; hence $Y'$, $X'$, and $f'$ also satisfy hypothesis $c)$ (\textbf{II},~7.3.8).
  \item By (\textbf{II},~7.3.8), the condition $c)$ of \sref{IV.15.7.5} may be replaced by the condition that $f'$ is \emph{proper}.
\end{enumerate}
\end{remarks}

\begin{env}[15.7.8]
\label{IV.15.7.8}
One says that a morphism of preschemes $f:X\to Y$ is \emph{submersive} if it is surjective and if the topology of $Y$ is equal to the quotient of the topology of $X$ by the equivalence relation defined by $f$.
One says that $f$ is \emph{universally submersive} if for every base change $Y'\to Y$, the morphism $f'=f_{(Y')}:X\times_Y Y'\to Y'$ is submersive.
Every surjective morphism \emph{universally open}, or \emph{universally closed}, or \emph{faithfully flat} and quasi-compact is universally submersive (2.3.12).
Given a morphism $f:X\to Y$, one says that a sub-set $Z$ of $X$ is \emph{submersive on $f(Z)$} if the topology of the sub-space $f(Z)$ of $Y$ is the quotient of the topology of the sub-space $Z$ of $X$ by the equivalence relation defined by $f|Z$.
One says that $Z$ is \emph{universally submersive on $f(Z)$} if, for every base change $g:Y'\to Y$, the image $Z'=p^{-1}(Z)$ of $Z$ by the projection $p:X\times_Y Y'\to Y'$ is a submersive set on $f'(Z')$.
One notes that if the morphism $f$ admits a $Y$-section $h:Y\to X$, then every set $Z$ containing $h(Y)$ is universally submersive on $Y$ (15.7.8).
\end{env}

\begin{proposition}[15.7.9]
\label{IV.15.7.9}
Let $Y$ be a locally Noetherian prescheme, $y$ a point of $Y$, $f:X\to Y$ a separated morphism of finite type.
Let $Z$ be a locally constructible part of $X$ having the following properties:
\begin{enumerate}[label=(\roman*)]
  \item For every generisation $y'$ of $y$, $Z_y=Z\cap f^{-1}(y')$ is a connected component of $f^{-1}(y')$ (4.5.2).
  \item $Z_y$ is proper on $\Spec(k(y))$.
  \item $Z$ is universally submersive on $f(Z)$ (15.7.8).
\end{enumerate}
Then $Z$ is proper at $y$ (in the sense of (15.7.1)), i.e.\ there exists a neighbourhood open $U$ of $y$ such that $Z\cap f^{-1}(U)$ is closed in $f^{-1}(U)$ and proper on $U$.
\end{proposition}

\begin{proof}
Using the method of (8.1.2,~$a$)) and (8.10.5,~(xii)), one reduces to proving the result when $Y=\Spec(A)$ is the spectrum of a \emph{local} Noetherian ring, and hypotheses (i), (ii), (iii) imply that $Z$ is proper on $Y$.

Note first that if $A'$ is a local Noetherian ring, $\varphi:A\to A'$ a local homomorphism, $Y'=\Spec(A')$ and $g:Y'\to Y$ the morphism corresponding to $\varphi$, the inverse image $Z'$ of $Z$ by the canonical projection $p:X\times_Y Y'\to X$ satisfies properties (i), (ii), (iii) of the statement, by (1.8.2) and (2.7.1,~(vii)); it therefore suffices to prove that $Z'$ is proper on $Y'$, provided that $A'=\widehat{A}$.

We use this remark by taking $A'=\widehat{A}$ ($\mathbf{0_I}$,~7.3.5), i.e.\ we may suppose that $A$ is \emph{complete}.
By (\textbf{III},~5.5.2), $X$ is the \emph{sum} of two sub-preschemes $X_0$, $X_1$ induced on open and closed parts of $X$, such that $X_0\cap f^{-1}(y)=Z_y$.
Set $Z_0=X_0\cap Z$, $Z_1=X_1\cap Z$; these sets form a partition of $Z$ into two open and closed parts in $Z$.
Since the $Z_y$ are connected, one necessarily has $Z_y\subset Z_0$ or $Z_y\subset Z_1$, hence $f(Z_0)\cap f(Z_1)=\emp$.
But since $Z_0$ and $Z_1$ are saturated for the equivalence relation defined by $f|Z$, it follows from (iii) that $f(Z_0)$ and $f(Z_1)$ are open in $f(Z)$.
Since $f(Z_0)$ contains $y$, and every neighbourhood of $y$ in $Y$ is equal to $Y$, one has $f(Z_0)=f(Z)$, and therefore $f(Z_1)=\emp$, i.e.\ $Z_1=\emp$ and $Z\subset X_0$.

One can now suppose that $f$ is \emph{proper}, and it remains to prove that a constructible part $Z$ of a prescheme $X$ proper over $Y$ satisfying only conditions (i) and (iii) is \emph{closed} in $X$.
By ($\mathbf{0_{III}}$,~9.2.5), it suffices to prove that for every $x'\in Z$ and every specialisation $x$ of $x'$ in $X$, one has $x\in Z$.
Let $A_1$ be a \emph{complete} discrete valuation ring, $u$ a morphism of $Y_1=\Spec(A_1)$ into $X$ such that, if $y_1$ and $y_1'$ are the closed point and generic point, $u(y_1)=x$ and $u(y_1')=x'$ (\textbf{II},~7.1.7); set $g=f\circ u:Y_1\to Y$ and $X_1=X\times_Y Y_1$, $f_1=f_{(Y_1)}$, where $p:X_1\to X$ is the canonical projection.
If $x_1=u_1(y_1)$, $x_1'=u_1(y_1')$, one has $p(x_1)=x$ and $p(x_1')=x'$, and $x_1'\in Z_1=p^{-1}(Z)$; since conditions (i) and (iii) are stable under all base changes, one reduces to the following situation: $Y$ is the spectrum of a complete discrete valuation ring, $X$ is proper over $Y$, there exists a $Y$-section $h$ of $X$ such that $h(y)=x$, $h(y')=x'$, and $x'\in Z$; one must prove that $x\in Z$.

Now, $Z_y$ is a connected component of $f^{-1}(y)$, proper on $\Spec(k(y))$ since $X$ is proper over $Y$.
Applying (\textbf{III},~5.5.2) again, there is an open and closed part $X'$ of $X$ such that $X'\cap f^{-1}(y)=Z_y$; since $Z_y$ is connected (replacing $X'$ by the connected component containing $Z_y$ if needed), one necessarily has $Z\subset X'$; since $h(Y)$ is connected and contains $x'\in Z\subset X'$, it is necessarily contained in $X'$, hence $x=h(y)\in X'\cap f^{-1}(y)=Z_y$.
\end{proof}

\begin{corollary}[15.7.10]
\label{IV.15.7.10}
Let $Y$ be a locally Noetherian prescheme, $f:X\to Y$ a separated morphism of finite type, universally submersive (15.7.8) and such that all the fibres $X_y=f^{-1}(y)$ are geometrically connected.
Then, if for every $y\in Y$ $X_y$ is proper on $k(y)$, $f$ is proper at the point $y$.
\end{corollary}

\begin{proof}
By \sref{IV.15.7.5}, one is reduced to the case where $Y=\Spec(\sh{O}_y)$, hence the hypotheses are stable by base change.
It suffices then to apply \sref{IV.15.7.9} to $Z=X$.
\end{proof}

\begin{corollary}[15.7.11]
\label{IV.15.7.11}
Let $f:X\to Y$ be a separated morphism of finite presentation; suppose that there exists a surjective morphism of finite presentation $g:Y'\to Y$, which is moreover either proper or flat, and such that if one sets $X'=X\times_Y Y'$, there exists a $Y'$-section of $X'$ (or even a $Y$-morphism $Y'\to X$ (\sref[I]{I.3.3.14})).
Suppose finally that the fibres $X_y=f^{-1}(y)$ are geometrically connected.
Then the set $U$ of the $y\in Y$ such that $X_y$ is proper on $k(y)$ is open in $Y$, and the restriction $f^{-1}(U)\to U$ of $f$ is proper.
\end{corollary}

\begin{proof}
\begin{lemma}[15.7.11.1]
\label{IV.15.7.11.1}
Let $f:X\to S$, $g:Y\to S$ be two morphisms, $p_1:X\times_S Y\to X$, $p_2:X\times_S Y\to Y$ the canonical projections.
If $s:S\to Y$ is a section of $g$, then $s'=(1,s\circ f)_S$ and $s':X\to X\times_S Y$, equipped with the morphisms $f:X\to S$ and $s':X\to X\times_S Y$, identifies with the fibre product of the $Y$-preschemes $S$ and $X\times_S Y$ for the morphisms $s:S\to Y$ and $p_2:X\times_S Y\to Y$.
\end{lemma}

This is a special case of (\sref[I]{I.3.3.11}), replacing the diagram by
\[
\xymatrix{
X \ar[r]^-{p_1} \ar[d] & X\times_S Y \ar[r]^-{s'} & X \ar[d]^{f}\\
S \ar[r]_-{s} & Y \ar[r]_-{g} & S
}
\]

Applying this lemma by replacing $S$, $Y$, $Y'$ by $Y$, $Y'$ respectively, one sees that $f=(f')_{(Y)}$ for the base change $s:Y\to Y'$, hence $f$ is affine since $f'$ is.
C.Q.F.D.
\end{proof}
